\subsection{Division on the common domains}

Vanishing on both center components gives a representation in the
generators $\bar a+\bar c$, $\nu$ and $\bar c(\bar B+\bar m)$.
We prove a finite-order integral form of this statement, compare these
generators with the passage coordinates, and apply it to transitions
on the common input domains.

\begin{proposition}[The center ideal and its local coordinates]
\label{rw:generator-division}
Put $\tau=\bar a+\bar c$, $\varsigma=\bar B+\bar m$ and
$\sigma=\bar c\varsigma$. Let $f$ be $C^{N+2}$ on a common
product coefficient box, with any section variables and $\bar B$
treated as passive variables. Suppose
\[
 f(0,0,\varsigma,0)=0,\qquad
 f(0,\bar c,0,0)=0,
\]
where the four arguments are $(\tau,\bar c,\varsigma,\nu)$.
Then
\begin{equation}\label{rw:centers:double-division}
 f=\tau A+\sigma B+\nu C,
\end{equation}
where the coefficients have bounded total order $N$ on a smaller
common box. Their bounds follow from the indicated derivatives of
$f$ on the larger box. A common bound for those derivatives by a
weight depending only on passive variables gives the same weight for
the divided coefficients.

For the finite three-saddle coordinates of
Lemma~\ref{rw:coordinates:finite-splittings}, there are constants
$0<c<C$ on a common smaller box such that
\begin{equation}\label{rw:centers:ideal-comparison}
 c\bigl(|\tau_1|+|\sigma_1|+|\nu|\bigr)
 \le |\xi_2|+|\kappa\xi_1|+|\nu|
 \le C\bigl(|\tau_1|+|\sigma_1|+|\nu|\bigr).
\end{equation}
Near a single component, the joint and weak coordinate blocks instead
give, respectively,
\begin{equation}\label{rw:centers:single-ideal-comparison}
\begin{aligned}
 |\tau_1|+|\sigma_1|+|\nu|
   &\asymp |\tau|+|\zeta|+|\nu|,\\
 |\tau_1|+|\sigma_1|+|\nu|
   &\asymp |d_s|+|\varsigma|+|\nu|.
\end{aligned}
\end{equation}
These comparisons concern coefficient functions on their inverse
boxes. Dividing transition expressions additionally uses the common
input domains constructed above.
\end{proposition}

\begin{proof}
First subtract $f(0,\bar c,\varsigma,0)$ by integrating in
$(\tau,\nu)$ along its straight segment:
\[
 f-f(0,\bar c,\varsigma,0)
 =\tau\int_0^1f_\tau(s\tau,\bar c,\varsigma,s\nu)\,\dd s
  +\nu\int_0^1f_\nu(s\tau,\bar c,\varsigma,s\nu)\,\dd s.
\]
The remaining function vanishes on both axes of its
$(\bar c,\varsigma)$ plane. Two applications of the fundamental
theorem give
\[
 f(0,\bar c,\varsigma,0)
 =\bar c\varsigma\int_0^1\!\int_0^1
   f_{\bar c\varsigma}(0,s\bar c,t\varsigma,0)\,\dd s\,\dd t.
\]
The segments lie in the prescribed box. Total-order differentiation
spends at most two coefficient derivatives and commutes with Euler
scale differentiation, proving the formula and its weighted version.

By \eqref{rw:coordinates:finite-center-axes}, apply this formula to
$\xi_2$ and $\kappa\xi_1$ in the original coefficient box for the
upper bound in \eqref{rw:centers:ideal-comparison}. Conversely the
same formula in the inverse product box
$(\xi_1,\xi_2,\kappa,\nu)$ gives, for $f=\tau_1,\sigma_1$,
\[
 f=\xi_2 A_f+\kappa\xi_1 B_f+\nu C_f,
\]
with bounded coefficients, proving the reverse bound.

In the joint block $\bar c$ has a positive lower bound, so
$|\sigma_1|\asymp|\varsigma|$. The exact inverse coordinate
map of Proposition~\ref{rw:center-coordinates} sends
$\tau=\zeta=\nu=0$ to
$\tau_1=\varsigma=\nu=0$. Integrating in these independent
transverse variables in each direction gives the first comparison
in \eqref{rw:centers:single-ideal-comparison}. In the weak block
the same argument uses $(d_s,\varsigma,\nu)$ and its inverse;
$\bar c$ again is a positive unit. These are single-component
divisions, each spending only one derivative.
\end{proof}

For transition division, retain larger signed intervals for the regular
outer legs and lower power inputs. Reduce the coefficient box so all
$G_\pm$ images stay there. Proposition~\ref{rw:lower-preparation}
gives the absolute zero extension at nonpositive deficits; extend the
relative coefficients by the same finite linear rule on both reflected
sides. Lower-deficit jets are ordinary. After singular $t^r$ composition,
the upper-input jets are Euler jets $t\partial_t$.

The following extension preserves the mixed total order.
Fix distinct positive numbers $a_0,\ldots,a_N$ and solve the
Vandermonde system
$\sum_{j=0}^N c_j(-a_j)^k=1$ for $0\le k\le N$.
For $x<0$ near zero set
\[
 (\mathcal Ef)(x;\mathbf b)=
       \sum_{j=0}^N c_j f(-a_jx;\mathbf b).
\]
Choose the negative interval so all $-a_jx$ lie in the positive
interval, with a cutoff farther from zero. The moments match input
jets through order $N$; the formula commutes with coefficient and
Euler derivatives and bounds mixed total order $N$ by that same
order. Constants depend only on the fixed order and interval.

The triangle identity holds on $u_0=z_0=\nu=0$, where every small
positive upper input has positive lower deficits. Reflection and the
common extension preserve the reversible identity at signed deficits.
Lemma~\ref{rw:coordinates:two-center-upper} preserves both in its
orbital remainder. Hence the integral paths
\[
 (s\xi_1,0,t\kappa,0),\qquad 0\le s,t\le1,
\]
and those removing $\xi_2,\nu$ lie in the auxiliary domains,
with no restriction to the physical max-norm face.

For the contact or weak comparison center, if
$[\mathbf b_c,\mathbf b]$ lies in the larger construction domain
and $f(p;\mathbf b_c)=0$ for every original input, then
\begin{equation}\label{rw:centers:one-center-taylor}
 f(p;\mathbf b)=
 \int_0^1D_{\mathbf b}f
       (p;\mathbf b_c+s(\mathbf b-\mathbf b_c))
       (\mathbf b-\mathbf b_c)\,\dd s.
\end{equation}
Propositions~\ref{rw:domains:contact} and \ref{rw:domains:weak}
bound the path length by the transverse size on the whole original
interval. Second-order remainders use second derivatives on these
same segments. The center choice is held fixed for the fibre.

\begin{lemma}[Relative forcing from a closing zero]
\label{rw:centers:relative-forcing}
On one fixed coefficient fibre let the central closing function be
\[
 F=Z-\rho\nu H,\qquad \rho=\sqrt\epsilon,
 \qquad H\ge h_*>0,
\]
where the reference gap $H$ is independent of the section input.
Suppose division on one of the common domains above gives
\begin{equation}\label{rw:centers:absolute-forcing-input}
 |Z|\le C(I+|\nu|)\epsilon^{-M}e^{-S_*/\epsilon},
 \qquad S_*>0,
\end{equation}
with a uniform positive action lower bound and an input-independent
nonnegative size $I$ of the relevant center generators. At any closing
zero one has
\begin{equation}\label{rw:centers:relative-forcing-output}
 |\nu|\le C' I\epsilon^{-M-1/2}e^{-S_*/\epsilon}.
\end{equation}
This is a constraint on the coefficients and hence holds on the whole
input fibre. If $I=0$, it forces $\nu=0$; where the maps have the
center identity there is no isolated closing zero.
\end{lemma}

\begin{proof}
At a zero of $F$, the left side of
\eqref{rw:centers:absolute-forcing-input} is
$\rho|\nu|H\ge\sqrt\epsilon h_*|\nu|$. Divide by this lower
factor. Since $S_*>0$ is fixed, the coefficient of $|\nu|$ on the
right tends uniformly to zero. For a small enough common scale it
is less than $1/2$ and can be absorbed. This proves
\eqref{rw:centers:relative-forcing-output}, including $I=0$.
\end{proof}

In particular, the two-center division preserves an absolute weight
$C\epsilon^{-M}e^{-S_*/\epsilon}$ with
$I=|\tau_1|+|\sigma_1|$. Section~4 supplies it when the retained
lower segment has a uniform positive action margin; that margin is
checked in each application.

The ideal comparisons transfer forcing to local generators:
$|\nu|\le\eta(I_{\rm global}+|\nu|)$ and
$I_{\rm global}\le C(I_{\rm local}+|\nu|)$ imply
$|\nu|\le2C\eta I_{\rm local}$ for small $\eta$.
The relation requires a closing zero.

On a block where $|\bar B+\bar m|$ has a positive lower bound,
$\bar c=\sigma_1/(\bar B+\bar m)$ and
$\bar a=\tau_1-\bar c$ also give a bounded comparison with
$(\bar a,\bar c)$. Near a finite upper saddle with regular lower
legs, $\bar a$ and $\kappa$ may replace these variables because
$\partial_{\bar c}\kappa=-1/\sqrt{\bar m}+O(\epsilon)$ there.
At a weak or a second-scale boundary the relevant generator retains
its small scale as a weight. We specify those weighted coordinates
when constructing the corresponding passages.
